\documentclass[11pt]{article}
\usepackage[margin=1in]{geometry}
\usepackage{amsmath, amssymb}
\usepackage{bm}
\usepackage{hyperref}
\usepackage{enumitem}
\usepackage{xcolor}

\title{My tryst with consistent stochastic modeling}
\author{Codex guided by Poovathingal}
\date{\today}

\begin{document}
\maketitle

\section{Big picture goal}
At a high-level, given a relationship between material properties and structure descriptors, and a way to sample the
structure descriptors, the problem is solved. The trivial solution to this problem are uniform bulk quantities. Given a single porosity, 
and analytical relationships between porosity and conductivity/permeability, obtaining the field of primitive variables is trivial. 


However, analytical expressions do not apply to our materials of interest. While higher order geometric descriptors such as fabric tensors could be used to generalize the 
property-structure relations, currently, based on analysis of permeability data, we can develop polynomial relationships that connect material property to porosity and a length scale. 

For stochastic modeling, the problem then reduces down to: {\it How do we sample porosity?} We will begin with generating a continuous co-variance function for porosity from the XRCT data. 

\section{Purpose}
This note describes the current computational strategy implemented in the
\texttt{porosity\_code} directory. The objective is to:
\begin{enumerate}[label=\arabic*.]
\item compute porosity from binarized TIFF stacks,
\item construct a spatial covariance description of the porosity field,
\item fit a parametric covariance model, and
\item generate continuous porosity realizations on a user-specified simulation grid.
\end{enumerate}

The writeup is organized around the same sequence used in the code. It is meant
to document both the mathematical intent and the practical approximations that
are currently being used.

\section{Block-Scale Permeability Interpretation}
First, let's talk about material properties. 


The permeability model
\[
K = K(\phi,L)
\]
should be interpreted at the block scale, not at the voxel scale, when the
training permeability data was obtained by cutting chunks of different lengths
from the original TIFF stacks.

In that case, the training data is naturally understood as relating:
\begin{itemize}
\item the porosity of a finite chunk,
\item the physical chunk length \(L\),
\item the permeability assigned to that chunk.
\end{itemize}

Thus the intended use of the permeability model is:
\[
K_B = K(\phi_B, L_B),
\]
where:
\begin{itemize}
\item \(B\) denotes a block or chunk,
\item \(\phi_B\) is the porosity averaged over that block,
\item \(L_B\) is the physical block length,
\item \(K_B\) is the permeability associated with that block.
\end{itemize}

After generating a continuous porosity field from the covariance model, the next
step is therefore to partition the realization into blocks, compute the average
porosity in each block, and then evaluate the permeability model on each block.

If the porosity realization is denoted by \(\phi(\bm{x})\), then for each block
\(B_j\) one computes:
\[
\phi_j
=
\frac{1}{|B_j|}
\int_{B_j}\phi(\bm{x})\,d\bm{x},
\]
or its discrete voxel-average equivalent. The permeability assigned to that block
is then:
\[
K_j = K(\phi_j, L_j).
\]

This produces a blockwise permeability field rather than a voxelwise
permeability field. The stochastic flow workflow is therefore:
\begin{enumerate}[label=\arabic*.]
\item generate a porosity realization from the covariance model,
\item divide the realization into blocks,
\item compute average porosity in each block,
\item evaluate \(K(\phi,L)\) for each block,
\item assemble the resulting blockwise permeability field,
\item solve flow through that heterogeneous medium.
\end{enumerate}

In this interpretation, the permeability model acts as a scale-dependent closure
that converts the covariance-generated porosity field into a permeability field
at the same physical scale used in the chunk-based permeability data.

\section{Input Data and Basic Assumptions}
The code assumes that each TIFF stack is already binarized. No grayscale
thresholding is performed in the porosity workflow. Instead, the stack is read
as a binary pore-solid volume.

Let the voxelized structure be denoted by
\[
I(\bm{x}) \in \{0,1\},
\]
where \(\bm{x} = (x,y,z)\) is a voxel index. The current implementation treats:
\begin{itemize}
\item \(I(\bm{x}) = 0\) as solid,
\item \(I(\bm{x}) = 1\) or any nonzero value as void.
\end{itemize}

Thus the binary field is interpreted directly as a pore-indicator field.

Two geometric quantities are supplied by the input file:
\begin{itemize}
\item \texttt{voxel\_resolution}: the physical size of one voxel,
\item \texttt{filter\_resolution}: the physical size of the local averaging
window used to construct the porosity field.
\end{itemize}

\section{Global Porosity Computation}
\subsection{Definition}
The global porosity is the fraction of void voxels in the entire stack:
\[
\phi_{\mathrm{global}}
=
\frac{1}{N}
\sum_{\bm{x}} I(\bm{x}),
\]
where \(N\) is the total number of voxels.

Since the input stack is already binary, this computation is straightforward:
\begin{itemize}
\item count all voxels,
\item count all nonzero voxels,
\item divide the second count by the first.
\end{itemize}

\subsection{Interpretation}
This quantity gives a single scalar summary of the sample. It does not provide
any spatial information. Two structures may have the same global porosity while
having very different spatial organization.

\section{Local Porosity Field}
\subsection{Motivation}
For covariance analysis and realization generation, the code does not work
directly with the raw binary field. Instead, it first constructs a continuous
local porosity field:
\[
\phi(\bm{x}),
\]
which measures the local void fraction in a neighborhood around \(\bm{x}\).

\subsection{Window-Based Averaging}
Let \(L\) denote the physical filter size. In the input file, this is
\texttt{filter\_resolution}. Let \(h\) denote the physical voxel size, given by
\texttt{voxel\_resolution}. The code converts the physical filter size into a
window size in voxels.

For a local window \(W_L(\bm{x})\) centered around \(\bm{x}\), the local
porosity is defined as:
\[
\phi(\bm{x})
=
\frac{1}{|W_L(\bm{x})|}
\sum_{\bm{u} \in W_L(\bm{x})} I(\bm{u}).
\]

This produces a continuous-valued field satisfying, ideally,
\[
0 \le \phi(\bm{x}) \le 1.
\]

\subsection{Role of the Filter Size}
The filter size controls the scale at which porosity fluctuations are observed:
\begin{itemize}
\item a small filter produces a more rapidly varying field,
\item a larger filter smooths local fluctuations and emphasizes larger-scale
structure.
\end{itemize}

Therefore, the covariance derived from the local porosity field depends
explicitly on the chosen filter size.

\section{Covariance of the Porosity Field}
\subsection{Mean and Variance}
Once the local porosity field \(\phi(\bm{x})\) is available, the code computes
its spatial mean:
\[
\mu = \langle \phi(\bm{x}) \rangle,
\]
and spatial variance:
\[
\sigma^2
=
\left\langle \left(\phi(\bm{x})-\mu\right)^2 \right\rangle.
\]

These are empirical averages over the finite voxel domain.

\subsection{Finite-Sample Spatial Covariance}
The covariance at lag \(\bm{r}\) is estimated by:
\[
C(\bm{r})
=
\frac{1}{N(\bm{r})}
\sum_{\bm{x}\in\Omega_{\bm{r}}}
\left(\phi(\bm{x})-\mu\right)
\left(\phi(\bm{x}+\bm{r})-\mu\right),
\]
where:
\begin{itemize}
\item \(\Omega_{\bm{r}}\) is the set of valid positions such that both
\(\bm{x}\) and \(\bm{x}+\bm{r}\) lie inside the domain,
\item \(N(\bm{r})\) is the number of valid overlapping pairs at lag
\(\bm{r}\).
\end{itemize}

This is a finite-sample covariance estimate. It is not an FFT-based covariance
and it does not use periodic boundary conditions.

\subsection{Normalized Covariance Used in the Code}
The lag covariance implementation stores a variance-normalized version:
\[
\rho(\bm{r}) = \frac{C(\bm{r})}{\sigma^2}.
\]
Thus:
\[
\rho(\bm{0}) \approx 1
\]
when the variance is nonzero.

This normalized quantity behaves like a spatial autocorrelation function. It is
useful for comparing the covariance shape independent of variance scale.

\section{3DLAG Covariance Construction}
\subsection{Motivation}
To support later realization generation, the code computes a full nonnegative
lag grid in the \texttt{3DLAG} mode:
\[
(r_x, r_y, r_z),
\qquad
0 \le r_x < N_x,\;
0 \le r_y < N_y,\;
0 \le r_z < N_z.
\]

Each grid point corresponds to one lag vector in voxel units, later converted to
physical units using \texttt{filter\_resolution}.

\subsection{Output Interpretation}
The 3DLAG table stores, for each lag:
\begin{itemize}
\item physical lag in \(x\), \(y\), and \(z\),
\item number of valid sample pairs,
\item mean porosity \(\mu\),
\item normalized covariance \(\rho(\bm{r})\).
\end{itemize}

The table therefore contains the information needed to reconstruct the covariance
field provided the variance \(\sigma^2\) is also known.

\section{Covariance Reconstruction}
\subsection{Step 1: Reconstruct the Covariance Field}
The preparation pipeline begins by reading the 3DLAG table and reconstructing the
unnormalized covariance:
\[
C(\bm{r}) = \sigma^2 \rho(\bm{r}).
\]

The mean \(\mu\) is not used in this reconstruction step because the
mean-centering has already been applied in the original covariance calculation.
However, \(\mu\) is carried forward because it is needed later when generating
continuous porosity realizations.

\subsection{Step 2: Symmetrize the Covariance}
The 3DLAG table contains only the nonnegative lag octant. A full signed lag
covariance field is needed for later spectral or large-domain operations.

The current code extends the nonnegative octant by reflection:
\[
C(r_x,r_y,r_z) = C(|r_x|,|r_y|,|r_z|).
\]

This is stronger than the stationarity identity
\[
C(\bm{r}) = C(-\bm{r}),
\]
because it also assumes reflection symmetry independently in each coordinate
direction. This is a practical approximation used to build a full signed grid
from the available data.

\section{Covariance Quality Checks}
\subsection{Tail Noise Analysis}
Finite-sample covariance estimates become less reliable at large lags because the
number of overlapping voxel pairs decreases as lag increases. The code therefore
analyzes the tail of the covariance field.

The current strategy is:
\begin{enumerate}[label=\arabic*.]
\item identify the outer portion of the lag domain,
\item bin tail samples by lag radius,
\item compute fluctuation statistics in each bin,
\item compare fluctuation size against local covariance magnitude and sample
count,
\item report whether the tail appears noisy.
\end{enumerate}

This does not prove validity or invalidity. It is a diagnostic intended to flag
lags that may be too noisy to trust strongly during model fitting.

\subsection{Negative Covariance Lobes}
The code also checks whether the empirical covariance becomes negative at some
lags. A negative covariance means that porosity fluctuations at the two points
tend to move in opposite directions relative to the mean.

This can be physically meaningful, for example when there is alternation or
layering at a characteristic spacing. However, a simple monotone exponential
covariance model cannot reproduce such oscillatory behavior. Therefore, the
presence of strong negative lobes is a signal that the exponential model may be
too simple.

\subsection{Anisotropy Analysis}
The code additionally examines directional behavior in order to assess whether the
covariance decay differs along \(x\), \(y\), and \(z\). This is useful because a
fiber-based material can be anisotropic.

The current first-pass realization model is isotropic, so the anisotropy report
acts as a warning mechanism. If the directional decay is significantly different,
an anisotropic covariance model would be more appropriate.

\section{Parametric Covariance Fitting}
\subsection{Why a Parametric Model is Needed}
The measured covariance is available only on a finite lag domain. If the desired
simulation domain is larger than the measured sample, covariance values are
needed at larger separations than those directly observed.

Rather than copying or periodically repeating the empirical covariance, the code
fits a smooth parametric model. The fitted model then serves as a covariance law
that can be evaluated on any desired simulation grid.

\textcolor{red}{Based on testing by Poovathingal, it is clear that the isotropic approach is not good. We need to consider other models.}

\subsection{Current Model: Isotropic Exponential Covariance}
The current fit uses an isotropic exponential form:
\[
C(r) = s \exp\left(-\frac{r}{\ell}\right),
\]
where:
\begin{itemize}
\item \(s\) is the sill,
\item \(\ell\) is the correlation length,
\item \(r = \|\bm{r}\|\) is the Euclidean lag magnitude.
\end{itemize}

This model is:
\begin{itemize}
\item positive at all lags,
\item monotone decreasing,
\item stationary and isotropic,
\item mathematically convenient for random field generation.
\end{itemize}

\subsection{Fitting Strategy}
The code currently fits the model using a log-linear least-squares approach over
strictly positive empirical covariance values:
\[
\log C(r) = \log s - \frac{r}{\ell}.
\]

Thus a straight-line fit is performed in \((r,\log C)\) space, yielding:
\[
\text{slope} = -\frac{1}{\ell},
\qquad
\text{intercept} = \log s.
\]

The current implementation therefore:
\begin{itemize}
\item excludes nonpositive empirical covariance values,
\item assumes the covariance is well-described by a single isotropic decay scale,
\item produces a first-pass fit rather than a fully general model.
\end{itemize}

\subsection{Fit Validation}
After fitting, the code validates the result by checking:
\begin{itemize}
\item whether the sill is positive,
\item whether the correlation length is positive,
\item whether enough points were used in the fit,
\item whether the fitted sill approximately matches the zero-lag empirical
covariance,
\item whether the RMS fitting error is reasonably small relative to the
zero-lag covariance scale.
\end{itemize}

This does not guarantee that the model is the best possible representation, but
it helps reject clearly unphysical or numerically unstable fits.

\section{Covariance Evaluation on the Simulation Grid}
\subsection{Need for a Larger Simulation Grid}
The simulation domain may be larger than the domain from which covariance was
measured. Therefore, the fitted model must be evaluated on the user-specified
realization grid:
\[
N_x \times N_y \times N_z,
\]
with physical dimensions:
\[
L_x \times L_y \times L_z.
\]

The code reads these from:
\begin{itemize}
\item \texttt{POROSITY\_REALIZATION\_GRID\_RES = \(N_x\ N_y\ N_z\)},
\item \texttt{POROSITY\_REALIZATION\_GRID\_DIM = \(L_x\ L_y\ L_z\)}.
\end{itemize}

\subsection{Grid Spacing}
The physical realization-grid spacings are:
\[
\Delta x = \frac{L_x}{N_x-1},
\qquad
\Delta y = \frac{L_y}{N_y-1},
\qquad
\Delta z = \frac{L_z}{N_z-1},
\]
for grids with more than one point in a direction.

The fitted covariance model is then evaluated over the simulation grid to build a
discrete covariance representation for realization generation.

\section{Continuous Covariance Realization Generation}
\subsection{Goal}
The goal is to generate many continuous porosity fields
\[
\phi^{(k)}(\bm{x}), \qquad k=1,2,\dots,N_{\mathrm{realizations}},
\]
that share the same target mean and covariance model while differing by random
draw.

\subsection{Why Gaussian Random Fields}
The current generation method uses a Gaussian random field. This is a practical
choice because, for a Gaussian field, the mean and covariance fully specify the
second-order statistics.

Thus once a fitted covariance model \(C(\bm{r})\) and a mean porosity \(\mu\)
are available, one can generate a zero-mean Gaussian field \(Z(\bm{x})\) with
the target covariance and then shift it by the mean:
\[
\phi(\bm{x}) = \mu + Z(\bm{x}).
\]

\subsection{Current Numerical Method}
The current implementation uses an explicit covariance matrix and Cholesky
factorization:
\begin{enumerate}[label=\arabic*.]
\item build the full covariance matrix for all realization-grid points,
\item compute its Cholesky factorization,
\item sample a vector of independent standard normal variables,
\item multiply by the Cholesky factor,
\item add the mean porosity.
\end{enumerate}

Formally, if
\[
\bm{z} \sim \mathcal{N}(\bm{0},\bm{I}),
\]
and if
\[
\bm{K} = \bm{L}\bm{L}^T
\]
is the Cholesky factorization of the covariance matrix \(\bm{K}\), then the
generated field vector is:
\[
\bm{\phi} = \mu \bm{1} + \bm{L}\bm{z}.
\]

\subsection{Current Practical Limitation}
This approach is exact for small to modest grids, but it scales poorly:
\begin{itemize}
\item memory grows as the square of the number of grid points,
\item factorization time grows cubically.
\end{itemize}

For that reason, the current code includes a hard point-count guard. Larger
realization grids would eventually need a more scalable FFT-based or
circulant-embedding-based generator.

\subsection{Generated Outputs}
The code writes realizations into a folder named
\texttt{GENERATED\_POROSITY\_FIELDS}. For each realization, two file types are
written:
\begin{itemize}
\item a TIFF stack for direct voxel-based storage,
\item a VTK file for visualization in ParaView.
\end{itemize}

The VTK output stores the continuous porosity field on the realization grid using
the user-specified physical dimensions and grid spacing.

\section{Overall Workflow Summary}
The full strategy implemented in the code can be summarized as follows:

\subsection*{Stage A: Porosity Computation}
\begin{enumerate}[label=\arabic*.]
\item Read a binarized TIFF stack.
\item Interpret zero voxels as solid and nonzero voxels as void.
\item Compute global porosity.
\item Compute a local porosity field using a window of physical size
\texttt{filter\_resolution}.
\end{enumerate}

\subsection*{Stage B: Covariance Construction}
\begin{enumerate}[label=\arabic*.]
\item Compute the mean and variance of the local porosity field.
\item Estimate finite-sample spatial covariance over valid overlapping pairs.
\item Store the normalized covariance in \texttt{3DLAG} mode.
\item Write tabular and ParaView-readable covariance outputs.
\end{enumerate}

\subsection*{Stage C: Covariance Preparation}
\begin{enumerate}[label=\arabic*.]
\item Read the 3DLAG table.
\item Reconstruct \(C(\bm{r}) = \sigma^2 \rho(\bm{r})\).
\item Symmetrize the covariance over a signed lag grid.
\item Analyze tail noise, negative lobes, and anisotropy.
\item Fit an isotropic exponential covariance model.
\item Validate the fitted model.
\item Evaluate the fitted covariance on the user-specified realization grid.
\end{enumerate}

\subsection*{Stage D: Covariance-Based Realization Generation}
\begin{enumerate}[label=\arabic*.]
\item Build the realization-grid covariance matrix from the fitted model.
\item Factorize the matrix by Cholesky decomposition.
\item Draw Gaussian random vectors.
\item Generate continuous porosity realizations by applying the covariance
factor and adding the target mean.
\item Write each realization as a TIFF stack and a VTK file.
\end{enumerate}

\section{Current Approximations and Caveats}
The present implementation is a good first end-to-end framework, but several
approximations should be kept in mind:
\begin{itemize}
\item The covariance is estimated from a finite sample, so large-lag values are
noisier.
\item The nonnegative 3DLAG octant is extended by reflection symmetry, which is
stronger than general stationarity.
\item The current parametric fit is isotropic and exponential, so it cannot
capture strong anisotropy or oscillatory negative covariance structure.
\item The Gaussian realization model produces a continuous field and does not
automatically enforce the physical bound \(0 \le \phi \le 1\).
\item The current Cholesky-based generator is appropriate only for relatively
small grids.
\end{itemize}

\section{Recommended Next Extensions}
For future refinement, the following extensions would be natural:
\begin{itemize}
\item anisotropic covariance fitting,
\item weighted covariance fitting using overlap-count information,
\item bounded transforms for porosity values,
\item FFT-based large-grid realization generation,
\item ensemble fitting across multiple TIFF samples.
\end{itemize}

\section{Closing Remark}
The overall strategy is to move from a measured binarized microstructure to a
continuous stochastic porosity model. The code does this by first converting the
raw stack into a local porosity field, then estimating its covariance, fitting a
smooth parametric model, and finally generating new continuous fields that match
the measured second-order statistics as closely as the current model allows.

\end{document}
